\section{Exact symbolic and finite controls}\label{app:certificates}

\subsection{Characteristic-zero smoothness transcript}

The coefficient-one recurrence used in
\cref{prop:smoothness} is replayed by the following complete Singular input.
The coefficient field is represented by adjoining \(r^2+r+1\); the
\texttt{dp} order is Singular's degree-reverse-lexicographic family.

\begin{verbatim}
ring R=0,(x0,x1,x2,x3,x4,x5,x6,x7,r),dp;
option(redSB);
poly Q=x0*x1+x1*x2+x2*x3+x3*x4+x4*x5+x5*x6
       +x6*x7+r*x7*x0;
ideal I=x0^2-x1-r*x7,
        x1^2-x0-x2,
        x2^2-x1-x3,
        x3^2-x2-x4,
        x4^2-x3-x5,
        x5^2-x4-x6,
        x6^2-x5-x7,
        x7^2-x6-r*x0,
        Q,
        r^2+r+1;
std(I);
\end{verbatim}

With Singular 4.2.1 and the reduced-standard-basis option, the exact output
is
\begin{verbatim}
_[1]=x7
_[2]=x6
_[3]=x5
_[4]=x4
_[5]=x3
_[6]=x2
_[7]=x1
_[8]=x0
_[9]=r^2+r+1
\end{verbatim}
No numerical tolerance or finite-prime extrapolation enters this
certificate.

\subsection{Bad-reduction control}

For \(p=181\), \(r=48\), one has
\[
 r^2+r+1\equiv0\pmod{181}.
\]
At
\[
 v=(9,158,158,9,104,128,171,153)
\]
the eight recurrence residuals are
\[
 (0,0,0,0,0,0,0,0)\quad\text{in }\mathbf F_{181},
\]
and direct substitution gives
\[
 \cQ(v)=\cC(v)=0.
\]
This negative control is theorem-bearing: it refutes all-split smoothness.
It does not contaminate the separate C48 curve factor.

\subsection{Chern coefficients}

The two Euler-characteristic inputs can be checked without a classification
table:
\[
 [H^6]\frac{(1+H)^8}{1+3H}=31,\qquad
 [H^5]\frac{(1+H)^8}{(1+2H)(1+3H)}=-27.                \tag{A.1}
\]
Multiplication by the degrees \(3\) and \(6\), followed by weak Lefschetz,
gives
\[
 \chi(S)=93,\quad b_6^{\mathrm{prim}}(S)=93-6-1=86,
\]
\[
 \chi(X)=-162,\quad b_5(X)=6-(-162)=168.
\]

\subsection{Finite factorization validation}

Finite controls are separated from the representation-theoretic proof.
For four split primes, exact extension-field counts reconstruct the full
degree-eight genus-four numerator by Newton identities:

\begin{center}
\begin{tabular}{@{}rrrrrr@{}}
\toprule
\(p\) & \(\#C(\mathbf F_p)\) & \(\#C(\mathbf F_{p^2})\) &
\(\#C(\mathbf F_{p^3})\) & \(\#C(\mathbf F_{p^4})\) &
\((a_+,a_-)\)\\
\midrule
7  & 12 & 66   & 372   & 2586   & \((-4,2)\)\\
13 & 18 & 270  & 2046  & 27414  & \((-1,-1)\)\\
19 & 24 & 474  & 6744  & 129930 & \((-4,2)\)\\
31 & 24 & 1050 & 29640 & 926970 & \((-4,8)\)\\
\bottomrule
\end{tabular}
\end{center}

In each row the reconstructed numerator equals
\[
 (1-a_+T+pT^2)^2(1-a_-T+pT^2)^2.
\]
An additional 21-prime degree-one ledger validates the trace splitting.
These ledgers confirm signs and normalizations; they do not replace
\cref{thm:jac-decomposition}.
